\documentclass[
	aspectratio=169,
	compress,
]{beamer}

\usetheme{Magdeburg}

\usepackage[T1]{fontenc}
\usepackage[utf8]{inputenc}

\usepackage{libertinus}
\usepackage[varqu,scaled=0.95]{inconsolata}
\usepackage{newtxmath}

\usepackage[english]{babel}

\usepackage{appendixnumberbeamer}
\usepackage{graphicx}
\usepackage[htt]{hyphenat}
\usepackage{listings}
\usepackage{microtype}
\usepackage{subcaption}
% Required for listings's upquote
\usepackage{textcomp}
\usepackage{upquote}
\usepackage{xcolor}

\usepackage[capitalise,noabbrev]{cleveref}

\usepackage[color=ovgu-orange]{todonotes}

\usepackage{lipsum}

\graphicspath{{./figures/}}

\lstset{
	basicstyle=\ttfamily\small,
	commentstyle=\color{ovgu-darkgray},
	keywordstyle=\color{ovgu-blue},
	numberstyle=\ttfamily\small\color{ovgu-darkgray},
	stringstyle=\color{ovgu-purple},
	rulecolor=\color{ovgu-lightgray},
	frame=single,
	numbers=left,
	language=C,
	breaklines=true,
	breakatwhitespace=true,
	postbreak=\hbox{$\hookrightarrow$ },
	showlines=true,
	showstringspaces=false,
	upquote=true,
	tabsize=4,
	gobble=4,
	captionpos=b,
	abovecaptionskip=\medskipamount,
}

\newcommand{\navframetitle}[1]{\frametitle{#1\hfill{\footnotesize\lastsection{}}}}

\title{Title}
\subtitle{Subtitle}

\author[Author]{
Author\\
{\footnotesize\href{mailto:author@ovgu.de}{\nolinkurl{author@ovgu.de}}}
}

\date{\today}

\institute{
Faculty of Computer Science\\
Otto von Guericke University Magdeburg
}

\titlegraphic{\vspace{0.9cm}\hfill\includegraphics[width=0.375\textwidth]{OVGU-INF}}

\AtBeginSection[]{
	\begin{frame}[plain,noframenumbering]
		\frametitle{Outline}

		\tableofcontents[currentsection]
	\end{frame}
}

\AtBeginSubsection[]{}

\begin{document}

\maketitle

\section{Introduction}
\label{sec:introduction}

\begin{frame}
	\navframetitle{Template}

	\begin{itemize}
		\item This presentation template is available at \url{https://github.com/parcio/templates} and consists of \cref{sec:introduction,sec:background,sec:evaluation,sec:conclusion}
	\end{itemize}
\end{frame}

\begin{frame}
	\navframetitle{Figures}

	\begin{figure}[ht]
		\centering
		\begin{subfigure}{0.45\textwidth}
			\centering
			\includegraphics[width=0.9\textwidth]{OVGU-INF}
			\caption{Left}
			\label{fig:left}
		\end{subfigure}
		\begin{subfigure}{0.45\textwidth}
			\centering
			\includegraphics[width=0.9\textwidth]{OVGU-INF}
			\caption{Right}
			\label{fig:right}
		\end{subfigure}
		\caption{Caption}
		\label{fig:both}
	\end{figure}

	\begin{itemize}
		\item You can refer to the subfigures (\cref{fig:left,fig:right}) or the figure (\cref{fig:both})
	\end{itemize}
\end{frame}

\section{Background}
\label{sec:background}

\begin{frame}
	\navframetitle{References}

	\begin{itemize}
		\item You can comfortably reference literature~\cite{DBLP:journals/superfri/DuweLMSF0B020}\footnote{This is a footnote.}
	\end{itemize}
\end{frame}

\begin{frame}
	\navframetitle{Tables}

	\begin{table}[ht]
		\centering
		\begin{tabular}{|l|c|r|}
			\hline
			\textbf{Header 1} & \textbf{Header 2} & \textbf{Header 3} \\
			\hline
			\hline
			Row 1 & Row 1 & Row 1 \\
			Row 2 & Row 2 & Row 2 \\
			\hline
		\end{tabular}
		\caption{Caption}
		\label{tab:table}
	\end{table}

	\begin{itemize}
		\item You can also refer to tables (\cref{tab:table})
	\end{itemize}
\end{frame}

\begin{frame}
	\navframetitle{Math}

	\[
		E = m c^2
	\]
\end{frame}

\section{Evaluation}
\label{sec:evaluation}

\begin{frame}[fragile]
	\navframetitle{Listings}

	\begin{lstlisting}[caption=Caption,label=lst:hello-world]
	printf("Hello world!\n");

	// Comment
	for (int i = 0; i < m; i++) {
		for (int j = 0; j < n; j++) {
			sum += 'a';
		}
	}
	\end{lstlisting}

	\begin{itemize}
		\item You can also refer to listings (\cref{lst:hello-world})
	\end{itemize}
\end{frame}

\begin{frame}[fragile]
	\navframetitle{Columns}

	\begin{columns}
		\begin{column}{0.5\textwidth}
			\begin{itemize}
				\item Slides can be split into columns
			\end{itemize}
		\end{column}
		\begin{column}{0.5\textwidth}
			\begin{lstlisting}[gobble=12]
			printf("Hello world!\n");

			// Comment
			for (int i = 0; i < m; i++) {
				for (int j = 0; j < n; j++) {
					sum += 'a';
				}
			}
			\end{lstlisting}
		\end{column}
	\end{columns}
\end{frame}

\section{Conclusion}
\label{sec:conclusion}

\begin{frame}
	\navframetitle{Todos}

	\todo[inline]{FIXME}

	\lipsum[1]
\end{frame}

\appendix

\AtBeginSection[]{}
\AtBeginSubsection[]{}

\section{References}

\begin{frame}[allowframebreaks]
	\frametitle{References}

	\bibliographystyle{apalike}
	\bibliography{presentation}
\end{frame}

\end{document}
